\documentclass[11pt,a4paper]{article}
% ============================================================
%  פרמבולה משותפת לחומרי הבחינה — מודולים 16–18
%  מהדרים עם:  lualatex <file>.tex
% ============================================================
\usepackage{fontspec}
\directlua{
  luaotfload.add_fallback("hbfb", {"DejaVu Serif:mode=node;", "DejaVu Sans:mode=node;"})
  luaotfload.add_fallback("ltfb", {"DejaVu Sans:mode=node;"})
}
\usepackage[bidi=basic]{babel}
\babelprovide[import, main]{hebrew}
\babelprovide[import]{english}
\babelfont{rm}[Ligatures=TeX, RawFeature={fallback=ltfb}]{DejaVu Serif}
\babelfont[hebrew]{rm}[Script=Hebrew, Ligatures=TeX, RawFeature={fallback=hbfb}]{Noto Serif Hebrew}
\babelfont{tt}[Scale=0.84, RawFeature={fallback=ltfb}]{DejaVu Sans Mono}

\usepackage[a4paper,top=2.1cm,bottom=2.0cm,left=1.9cm,right=1.9cm,headsep=0.5cm,footskip=1cm]{geometry}
\usepackage{booktabs, longtable, array, adjustbox, xcolor, enumitem, fancyhdr, titlesec,
            fancyvrb, needspace, multicol, lastpage, tikz, colortbl}
\usetikzlibrary{positioning, fit, backgrounds, matrix}

\definecolor{inkblue}{HTML}{14365A}
\definecolor{rulegray}{HTML}{8A97A3}
\definecolor{boxgray}{HTML}{F4F6F8}
\definecolor{codebg}{HTML}{F7F8FA}
\definecolor{warnred}{HTML}{9B2226}
\definecolor{okgreen}{HTML}{1B5E3A}
\definecolor{tipgold}{HTML}{8A6D1F}

\setlength{\parindent}{0pt}
\setlength{\parskip}{3pt plus 1pt}
\renewcommand{\arraystretch}{1.15}
\setlength{\columnsep}{18pt}
\setlength{\headheight}{13pt}

\pagestyle{fancy}\fancyhf{}
\renewcommand{\headrulewidth}{0.4pt}
\renewcommand{\footrulewidth}{0pt}
\fancyfoot[C]{\small עמוד \thepage{} מתוך \pageref{LastPage}}

% ---- כותרות ----
\titleformat{\section}{\needspace{3\baselineskip}\large\bfseries\color{inkblue}}{\thesection.\ }{0pt}{}
\titlespacing*{\section}{0pt}{15pt plus 3pt}{6pt}
\titleformat{\subsection}{\needspace{2.5\baselineskip}\normalsize\bfseries\color{inkblue}}{\thesubsection\ }{0pt}{}
\titlespacing*{\subsection}{0pt}{10pt plus 2pt}{4pt}

% ---- קו מפריד מתחת לכותרת חלק ----
\newcommand{\partheader}[3]{%
  \section*{#1}%
  \vspace{-4pt}{\color{rulegray}\hrule height 0.8pt}\vspace{4pt}%
  {\small\textbf{בסיס הנתונים:} #2 \hfill \textbf{ניקוד:} #3}\par\vspace{2pt}%
  {\color{rulegray}\hrule height 0.4pt}\vspace{6pt}}

\newcommand{\rulebar}{\vspace{-4pt}{\color{rulegray}\hrule height 0.8pt}\vspace{5pt}}

% ---- שאלה וסעיף ----
\newcounter{qnum}
\newcommand{\question}[2]{%
  \needspace{5\baselineskip}\refstepcounter{qnum}\par\vspace{8pt}%
  \noindent\textbf{\color{inkblue}שאלה \theqnum.}\ \textbf{#2}\hfill{\small(#1 נק')}\par\vspace{2pt}}
\newcommand{\subq}[2]{\needspace{3\baselineskip}\par\vspace{4pt}\noindent\textbf{#1.}~#2\par}
\newcommand{\subqp}[3]{\needspace{8\baselineskip}\par\vspace{5pt}\noindent%
  \textbf{#1.}~{\small\textbf{\color{inkblue}(#3 נק')}}~#2\par}
\newcommand{\pairitem}[1]{\needspace{2\baselineskip}\par\vspace{4pt}\noindent\hspace{1.4em}\textbf{#1}\par}

% ---- שורות תשובה מנוקדות ----
\newcommand{\answerlines}[1]{%
  \par\vspace{1pt}%
  \needspace{\dimexpr 6mm*#1 + 2mm\relax}%
  \begingroup\color{rulegray}%
  \count255=#1\relax
  \loop\ifnum\count255>0
    \noindent\babelsublr{\rule[-2pt]{0.6pt}{9pt}\rule[-2pt]{\dimexpr\linewidth-0.6pt}{0.3pt}}\par\vspace{5.4mm}%
    \advance\count255 by -1
  \repeat
  \endgroup\vspace{-2mm}}

% ---- שורת תשובה קצרה אחת ----
\newcommand{\blankline}[1]{\babelsublr{\rule[-2pt]{#1}{0.4pt}}}

% ---- קוד SQL, תמיד משמאל לימין ----
\DefineVerbatimEnvironment{sqlverb}{Verbatim}{fontsize=\small, frame=leftline, framerule=1.2pt,
  rulecolor=\color{rulegray}, framesep=7pt, xleftmargin=12pt, samepage=true}
\DefineVerbatimEnvironment{sqlbad}{Verbatim}{fontsize=\small, frame=leftline, framerule=1.2pt,
  rulecolor=\color{warnred}, framesep=7pt, xleftmargin=12pt, samepage=true}
\DefineVerbatimEnvironment{sqlgood}{Verbatim}{fontsize=\small, frame=leftline, framerule=1.2pt,
  rulecolor=\color{okgreen}, framesep=7pt, xleftmargin=12pt, samepage=true}

\newcommand{\lr}[1]{\babelsublr{\mbox{#1}}}
\newcommand{\sqlin}[1]{\babelsublr{\texttt{#1}}}

% ---- תיבה אפורה ----
\newsavebox{\panelbox}
\newenvironment{panel}{\par\addvspace{7pt}\begin{lrbox}{\panelbox}%
  \begin{minipage}{\dimexpr\linewidth-2\fboxsep-2\fboxrule}}%
  {\end{minipage}\end{lrbox}\noindent\fcolorbox{rulegray}{boxgray}{\usebox{\panelbox}}\par\addvspace{7pt}}

% ---- תיבה צבעונית עם כותרת ----
\newsavebox{\notebox}
\newenvironment{notepanel}[2]{\def\npcolor{#1}\par\addvspace{7pt}\begin{lrbox}{\notebox}%
  \begin{minipage}{\dimexpr\linewidth-2\fboxsep-2\fboxrule}%
  \textbf{\color{#1}#2}\par\vspace{2pt}}%
  {\end{minipage}\end{lrbox}\noindent\fcolorbox{\npcolor}{white}{\usebox{\notebox}}\par\addvspace{7pt}}

% ---- טבלת נתונים מלאה ----
\newcommand{\thd}[1]{\textbf{#1}}
\newenvironment{fulltable}[3]{%
  \needspace{5\baselineskip}\par\addvspace{7pt}%
  \noindent{\small הטבלה \babelsublr{\texttt{\textbf{#1}}} \textendash\ #2#3}\par\vspace{2pt}%
  \begin{otherlanguage}{english}\begin{adjustbox}{max width=\linewidth}}
  {\end{adjustbox}\end{otherlanguage}\par\addvspace{4pt}}

% ---- פלט של שאילתה ----
\newenvironment{outtable}[1]{%
  \needspace{4\baselineskip}\par\vspace{2pt}%
  \noindent{\small\textbf{הפלט}#1}\par\vspace{1pt}%
  \begin{otherlanguage}{english}\begin{adjustbox}{max width=\linewidth}}
  {\end{adjustbox}\end{otherlanguage}\par\addvspace{5pt}}

% ---- פלט ללא כותרת ----
\newenvironment{outplain}{%
  \par\vspace{2pt}\begin{otherlanguage}{english}\begin{adjustbox}{max width=\linewidth}}
  {\end{adjustbox}\end{otherlanguage}\par\addvspace{4pt}}

\newcommand{\nullv}{\textit{\small NULL}}
\providecommand{\tightlist}{\setlength{\itemsep}{0pt}\setlength{\parskip}{0pt}}


\fancyhead[R]{\small בסיסי נתונים ו\lr{-}\lr{SQL} \textbar{} בחינת אמצע}
\fancyhead[L]{\small \lr{\texttt{SELECT} · \texttt{WHERE} · \texttt{ORDER BY}}}

\begin{document}

\begin{center}
{\LARGE\bfseries\color{inkblue} בחינת אמצע בבסיסי נתונים}\\[3pt]
{\large טבלאות ומפתחות \textbullet{} \texttt{SELECT} \textbullet{} \texttt{WHERE} \textbullet{} \texttt{ORDER BY}}\\[6pt]
{\large\bfseries משך הבחינה: שעתיים \ \textbar\ \ \lr{100} נקודות}
\end{center}
\vspace{2pt}
{\color{inkblue}\hrule height 1.2pt}
\vspace{8pt}

\noindent שם המגיש/ה:~\rule{5.2cm}{0.4pt}\hfill כיתה:~\rule{2.2cm}{0.4pt}\hfill תאריך:~\rule{3cm}{0.4pt}
\vspace{8pt}

\begin{panel}
\textbf{הוראות}
\begin{itemize}[leftmargin=1.3em, itemsep=2pt, topsep=3pt]
  \item בבחינה \textbf{ארבע שאלות}, ועונים על \textbf{כולן}. הניקוד של כל סעיף רשום לצדו. בסוף יש סעיף \textbf{בונוס} של \lr{5} נקודות.
  \item \textbf{מותר להשתמש בחומר עזר} כתוב או מודפס \textendash{} דף העזר, המחברת וחומר הלימוד. \textbf{אסור} להשתמש במחשב או בטלפון.
  \item \textbf{כל השאילתות בבחינה הזאת הן על טבלה אחת.} אין צורך לצרף טבלאות.
  \item \textbf{שפת \lr{SQL} נכתבת משמאל לימין.} מתחילים לכתוב מהקצה ה\textbf{שמאלי} של השורה.
  \item ערך טקסט נכתב בגרש \textbf{בודד}: \sqlin{'דחוף'}. גרש \textbf{כפול} הוא שם של עמודה או כינוי: \sqlin{AS "מחיר"}.
  \item מסיימים כל שאילתה ב\textbf{נקודה-פסיק}.
  \item תאריך נכתב תמיד בצורה \sqlin{'YYYY-MM-DD'}, למשל \sqlin{'2026-03-01'}.
\end{itemize}
\end{panel}

\vspace{6pt}
\section*{בסיס הנתונים: מעבדת "דלתא-טק"}
\vspace{-5pt}{\color{rulegray}\hrule height 0.8pt}\vspace{6pt}

"דלתא-טק" היא מעבדה לתיקון מכשירים. לקוח מביא מכשיר, המעבדה פותחת \textbf{תיקון}, ומשייכת אותו ל\textbf{טכנאי} אחד. במהלך התיקון נכנסים למכשיר \textbf{חלפים}. לכל טכנאי יש \textbf{עמדת עבודה} אחת לכל היותר.

\vspace{3pt}
\noindent\textbf{הטבלאות:}
\begin{itemize}[leftmargin=1.3em, itemsep=1pt, topsep=2pt]
  \item \textbf{טבלאות ראשיות:} \lr{\texttt{Technicians}} (טכנאים) ו-\lr{\texttt{Repairs}} (תיקונים).
  \item \textbf{טבלאות עזר:} \lr{\texttt{Cities}} (ערים), \lr{\texttt{DevTypes}} (סוגי מכשירים) ו-\lr{\texttt{Parts}} (חלפים).
  \item \textbf{טבלאות נוספות:} \lr{\texttt{Benches}} (עמדות עבודה) ו-\lr{\texttt{PartUsage}} (אילו חלפים נכנסו לאיזה תיקון, ובאיזו כמות).
\end{itemize}

\vspace{2pt}
{\small\textbf{הערות לנתונים:} \lr{\texttt{HourWage}} הוא שכר הטכנאי לשעה בש"ח. \lr{\texttt{RepairPrice}} הוא מחיר התיקון בש"ח. \lr{\texttt{InDate}} הוא תאריך קבלת המכשיר, ו-\lr{\texttt{DoneDate}} הוא תאריך סיום התיקון \textendash{} \textbf{ערך \lr{\texttt{NULL}} פירושו שהתיקון עדיין פתוח}. לשני טכנאים לא רשום טלפון.}

\vspace{4pt}
\noindent\begin{minipage}[t]{0.48\linewidth}
\begin{fulltable}{Cities}{6 שורות}{}
\begin{tabular}{ll}
\toprule
\thd{CityCode} & \thd{CityName} \\
\midrule
1 & ירכא \\
2 & עכו \\
3 & כרמיאל \\
4 & נהריה \\
5 & שפרעם \\
6 & מעלות \\
\bottomrule
\end{tabular}
\end{fulltable}
\end{minipage}\hfill
\begin{minipage}[t]{0.48\linewidth}
\begin{fulltable}{Parts}{7 שורות}{}
\begin{tabular}{lll}
\toprule
\thd{PartCode} & \thd{PartName} & \thd{PartPrice} \\
\midrule
101 & מסך & 450 \\
102 & סוללה & 180 \\
103 & שקע טעינה & 90 \\
104 & רמקול & 70 \\
105 & מקלדת & 220 \\
106 & מצלמה & 160 \\
107 & כבל שטוח & 45 \\
\bottomrule
\end{tabular}
\end{fulltable}
\end{minipage}

\vspace{2pt}
\begin{fulltable}{DevTypes}{6 שורות}{}
\begin{tabular}{lll}
\toprule
\thd{DevCode} & \thd{DevName} & \thd{DevFamily} \\
\midrule
10 & טלפון חכם & סלולר \\
20 & מחשב נייד & מחשבים \\
30 & טאבלט & סלולר \\
40 & שעון חכם & לבישים \\
50 & אוזניות & לבישים \\
60 & קונסולת משחק & מחשבים \\
\bottomrule
\end{tabular}
\end{fulltable}

\begin{fulltable}{Technicians}{7 שורות}{}
\begin{tabular}{lllllll}
\toprule
\thd{TechCode} & \thd{TechFirst} & \thd{TechLast} & \thd{TechCity} & \thd{StartDate} & \thd{HourWage} & \thd{TechPhone} \\
\midrule
1 & נביל & סרחאן & 1 & 2019-03-01 & 95 & 050-7010101 \\
2 & רונית & בר-לב & 4 & 2021-09-15 & 80 & \nullv \\
3 & חוסאם & זיאד & 5 & 2020-01-12 & 88 & 054-7030303 \\
4 & דנה & כהן & 1 & 2023-06-01 & 72 & 052-7040404 \\
5 & סאמר & חלבי & 1 & 2018-11-20 & 105 & \nullv \\
6 & מאיה & אברהם & 6 & 2024-02-10 & 65 & 053-7060606 \\
7 & ליאן & קבלאן & 2 & 2022-05-05 & 80 & 050-7070707 \\
\bottomrule
\end{tabular}
\end{fulltable}

\begin{fulltable}{Repairs}{16 שורות}{}
\begin{tabular}{llllllll}
\toprule
\thd{RepairNum} & \thd{CustName} & \thd{RepairDev} & \thd{RepairTech} & \thd{InDate} & \thd{DoneDate} & \thd{RepairPrice} & \thd{Priority} \\
\midrule
1001 & אדם חלבי & 10 & 1 & 2026-01-12 & 2026-01-15 & 420 & רגיל \\
1002 & מאיה כהן & 20 & 3 & 2026-01-20 & 2026-02-02 & 980 & דחוף \\
1003 & כרים עזאם & 10 & 2 & 2026-02-03 & 2026-02-05 & 310 & רגיל \\
1004 & שירה לוי & 30 & 1 & 2026-02-14 & \nullv & 540 & דחוף \\
1005 & רמי חורי & 60 & 5 & 2026-03-02 & 2026-03-09 & 760 & רגיל \\
1006 & תמר אברהם & 40 & 4 & 2026-03-11 & 2026-03-12 & 190 & מיידי \\
1007 & ליאן קבלאן & 10 & 7 & 2026-03-25 & 2026-04-01 & 450 & רגיל \\
1008 & עומר ביטון & 20 & 3 & 2026-04-05 & \nullv & 1250 & דחוף \\
1009 & סאלי חסון & 50 & 6 & 2026-04-18 & 2026-04-19 & 120 & רגיל \\
1010 & יואב שושן & 10 & 1 & 2026-05-02 & 2026-05-06 & 380 & מיידי \\
1011 & נור מנסור & 30 & 2 & 2026-05-15 & 2026-05-20 & 610 & רגיל \\
1012 & אלון כהן & 20 & 5 & 2026-06-01 & 2026-06-14 & 1480 & דחוף \\
1013 & סלים דיאב & 10 & 7 & 2026-06-22 & 2026-06-24 & 290 & רגיל \\
1014 & הילה פרץ & 60 & 4 & 2026-07-03 & \nullv & 820 & רגיל \\
1015 & מוחמד דיב & 40 & 6 & 2026-07-19 & 2026-07-21 & 230 & מיידי \\
1016 & אורי שמש & 20 & 3 & 2026-08-04 & 2026-08-20 & 1100 & דחוף \\
\bottomrule
\end{tabular}
\end{fulltable}

\vspace{2pt}
\noindent\begin{minipage}[t]{0.40\linewidth}
\begin{fulltable}{Benches}{5 שורות}{}
\begin{tabular}{lll}
\toprule
\thd{BenchNum} & \thd{BenchTech} & \thd{BenchFloor} \\
\midrule
1 & 1 & 1 \\
2 & 3 & 1 \\
3 & 5 & 2 \\
4 & 7 & 2 \\
5 & 4 & 2 \\
\bottomrule
\end{tabular}
\end{fulltable}
\end{minipage}\hfill
\begin{minipage}[t]{0.56\linewidth}
\begin{fulltable}{PartUsage}{15 שורות}{}
\begin{tabular}{lll|lll}
\toprule
\thd{UseRepair} & \thd{UsePart} & \thd{UseQty} & \thd{UseRepair} & \thd{UsePart} & \thd{UseQty} \\
\midrule
1001 & 102 & 1 & 1008 & 105 & 1 \\
1002 & 105 & 1 & 1010 & 102 & 1 \\
1002 & 102 & 1 & 1011 & 101 & 1 \\
1003 & 101 & 1 & 1012 & 105 & 1 \\
1004 & 101 & 1 & 1012 & 102 & 1 \\
1005 & 104 & 2 & 1016 & 101 & 1 \\
1007 & 103 & 1 & 1016 & 106 & 1 \\
1008 & 101 & 1 & & & \\
\bottomrule
\end{tabular}
\end{fulltable}
\end{minipage}

\newpage
\question{20}{מבנה בסיס הנתונים}

\subqp{א}{כתבו את \textbf{המפתח הראשי} של כל אחת מהטבלאות הבאות. שימו לב: באחת מהן המפתח \textbf{מורכב} משני שדות.}{4}
\vspace{2pt}
\noindent\begin{tabular}{@{}p{0.22\linewidth}p{0.70\linewidth}@{}}
\lr{\texttt{Technicians}} & \blankline{7cm} \\[7pt]
\lr{\texttt{Repairs}} & \blankline{7cm} \\[7pt]
\lr{\texttt{Benches}} & \blankline{7cm} \\[7pt]
\lr{\texttt{PartUsage}} & \blankline{7cm} \\[5pt]
\end{tabular}

\subqp{ב}{בטבלת \lr{\texttt{Repairs}} יש \textbf{שני} מפתחות זרים. כתבו את שם כל אחד מהם, ואת \textbf{הטבלה והשדה} שאליהם הוא מצביע.}{4}
\vspace{2pt}
\noindent\begin{tabular}{@{}p{0.30\linewidth}p{0.30\linewidth}p{0.32\linewidth}@{}}
\toprule
\thd{המפתח הזר} & \thd{מצביע לטבלה} & \thd{לשדה} \\
\midrule
\blankline{3.4cm} & \blankline{3.4cm} & \blankline{3.4cm} \\[8pt]
\blankline{3.4cm} & \blankline{3.4cm} & \blankline{3.4cm} \\[4pt]
\bottomrule
\end{tabular}

\subqp{ג}{לפניכם שלושה זוגות של טבלאות. כתבו לכל זוג את \textbf{סוג הקשר} (\lr{1:1} או \lr{1:N} או \lr{M:N}), \textbf{ונמקו במשפט אחד}.}{6}
\vspace{1pt}
\pairitem{\lr{\texttt{Cities} \ \textendash\ \ \texttt{Technicians}}}
\answerlines{2}
\pairitem{\lr{\texttt{Repairs} \ \textendash\ \ \texttt{Parts}}}
\answerlines{2}
\pairitem{\lr{\texttt{Technicians} \ \textendash\ \ \texttt{Benches}}}
\answerlines{2}

\subqp{ד}{לפניכם תרשים \lr{DSD} חלקי. \textbf{ציירו חץ} מכל מפתח זר אל המפתח הראשי שהוא מצביע עליו (שלושה חצים).}{3}

\vspace{4pt}
\begin{otherlanguage}{english}
\begin{center}
\begin{tikzpicture}[font=\scriptsize\sffamily]
\tikzset{tblm/.style={matrix of nodes,
  nodes={anchor=west, inner sep=3pt, text width=26mm, align=left, font=\scriptsize\sffamily},
  row sep=-\pgflinewidth, draw=inkblue, line width=0.7pt, inner sep=0pt,
  row 1/.style={nodes={fill=inkblue, text=white, font=\scriptsize\sffamily\bfseries,
                       align=center, text width=26mm}}}}
\matrix (C) [tblm] { Cities \\ CityCode \\ CityName \\ };
\matrix (T) [tblm, right=22mm of C] { Technicians \\ TechCode \\ TechFirst \\ TechLast \\ TechCity \\ StartDate \\ HourWage \\ };
\matrix (R) [tblm, right=22mm of T] { Repairs \\ RepairNum \\ CustName \\ RepairDev \\ RepairTech \\ InDate \\ RepairPrice \\ };
\matrix (D) [tblm, below=9mm of C] { DevTypes \\ DevCode \\ DevName \\ DevFamily \\ };
\end{tikzpicture}
\end{center}
\end{otherlanguage}

\subqp{ה}{נכון או לא נכון? הקיפו, ו\textbf{הסבירו במשפט} כל קביעה שסימנתם "לא נכון".}{3}
\vspace{2pt}
\subq{1}{אפשר להוסיף לטבלת \lr{\texttt{Benches}} שורה נוספת שבה \sqlin{BenchTech = 1}. \hfill \textbf{נכון \ / \ לא נכון}}
\subq{2}{בטבלת \lr{\texttt{PartUsage}}, השדה \lr{\texttt{UseRepair}} לבדו יכול לשמש מפתח ראשי. \hfill \textbf{נכון \ / \ לא נכון}}
\subq{3}{יש עיר בטבלת \lr{\texttt{Cities}} שאין בה אף טכנאי. \hfill \textbf{נכון \ / \ לא נכון}}
\answerlines{3}

\vspace{10pt}
\question{25}{\lr{\texttt{SELECT}} \textendash{} מה מציגים}

\subqp{א}{הציגו מטבלת \lr{\texttt{Repairs}} את מספר התיקון, שם הלקוח ומחיר התיקון. תנו לשלוש העמודות כותרות בעברית: \textbf{מספר תיקון}, \textbf{שם הלקוח} ו-\textbf{מחיר}.}{5}
\answerlines{3}

\subqp{ב}{אילו רמות דחיפות קיימות בטבלת התיקונים? הציגו כל רמה \textbf{פעם אחת בלבד}.}{5}
\answerlines{2}

\subqp{ג}{הציגו את השם המלא של כל טכנאי \textendash{} שם פרטי, רווח, ושם משפחה \textendash{} בעמודה אחת, תחת הכותרת \lr{\texttt{FullName}}.}{5}
\answerlines{3}

\subqp{ד}{למחיר כל תיקון מתווסף מע"מ בשיעור \lr{18\%}. הציגו את מספר התיקון, את המחיר, ואת המחיר כולל מע"מ תחת הכותרת \lr{\texttt{TotalPrice}}.}{5}
\answerlines{3}

\subqp{ה}{השאילתה הבאה \textbf{אינה רצה}. יש בה \textbf{שלוש} טעויות. כתבו אותה מחדש, מתוקנת.}{5}

\begin{otherlanguage}{english}
\begin{sqlbad}
SELECT RepairNum, CustName, Priority,
FROM   Repairs
WHERE  Priority = "רגיל"
\end{sqlbad}
\end{otherlanguage}
\answerlines{4}

\newpage
\question{30}{\lr{\texttt{WHERE}} \textendash{} אילו שורות}

\subqp{א}{הציגו את מספר התיקון, שם הלקוח והמחיר, של כל התיקונים שמחירם \textbf{גבוה מ-\lr{700}} ש"ח.}{5}
\answerlines{3}

\subqp{ב}{הציגו את מספר התיקון, תאריך הקבלה והדחיפות, של התיקונים שהתקבלו במעבדה \textbf{בין \lr{1.3.2026} ל-\lr{31.5.2026}} (כולל). השתמשו ב-\lr{\texttt{BETWEEN}}.}{5}
\answerlines{3}

\subqp{ג}{הציגו את מספר התיקון, שם הלקוח והדחיפות, של התיקונים שדחיפותם \textbf{"דחוף" או "מיידי"}. השתמשו ב-\lr{\texttt{IN}}.}{5}
\answerlines{3}

\subqp{ד}{הציגו את מספר התיקון ואת שם הלקוח, של כל הלקוחות ש\textbf{שם המשפחה שלהם "כהן"}. שימו לב: בעמודה \lr{\texttt{CustName}} שמור השם המלא.}{5}
\answerlines{3}

\subqp{ה}{הציגו את מספר התיקון, שם הלקוח ותאריך הקבלה, של כל התיקונים ש\textbf{עדיין פתוחים}.}{5}
\answerlines{3}

\subqp{ו}{הציגו את שם המשפחה ואת השכר לשעה, של הטכנאים ש\textbf{גרים בירכא (קוד \lr{1}) או בעכו (קוד \lr{2})}, \textbf{ו}שהשכר שלהם לשעה הוא \textbf{\lr{80} ש"ח ומעלה}.}{5}

\vspace{-2pt}
{\small\textbf{שימו לב} \textendash{} בסעיף הזה יש \lr{\texttt{OR}} ו-\lr{\texttt{AND}} באותה שאילתה.}
\answerlines{3}

\vspace{10pt}
\question{25}{\lr{\texttt{ORDER BY}} \textendash{} באיזה סדר}

\subqp{א}{הציגו את מספר התיקון, הדחיפות והמחיר של \textbf{כל} התיקונים, ממוינים לפי \textbf{הדחיפות} בסדר עולה (לפי האלף-בית), ובתוך כל דחיפות \textendash{} לפי \textbf{המחיר מהגבוה לנמוך}.}{5}
\answerlines{3}

\subqp{ב}{הציגו את \textbf{שלושת התיקונים היקרים ביותר} מבין התיקונים ש\textbf{כבר הסתיימו}: מספר תיקון, שם לקוח ומחיר.}{5}
\answerlines{3}

\subqp{ג}{הציגו את שם המשפחה והשכר לשעה של כל הטכנאים, ממוינים לפי \textbf{השכר מהגבוה לנמוך}. לשני טכנאים אותו שכר \textendash{} ביניהם מיינו לפי \textbf{שם המשפחה} בסדר האלף-בית.}{5}
\answerlines{3}

\subqp{ד}{הציגו את שם החלף ואת מחירו, לחלפים שמחירם \textbf{בין \lr{100} ל-\lr{300}} ש"ח, \textbf{מהיקר לזול}.}{5}
\answerlines{3}

\subqp{ה}{לפניכם שאילתה. ענו על שתי השאלות שאחריה.}{5}

\begin{otherlanguage}{english}
\begin{sqlverb}
SELECT   RepairNum, DoneDate
FROM     Repairs
ORDER BY DoneDate;
\end{sqlverb}
\end{otherlanguage}

\subq{1}{\textbf{כמה שורות} תחזיר השאילתה?}
\answerlines{1}
\subq{2}{\textbf{אילו שורות יופיעו ראשונות, ומדוע?} כתבו את מספרי התיקונים.}
\answerlines{3}

\vspace{6pt}
{\color{rulegray}\hrule height 0.8pt}
\vspace{4pt}
\question{5}{בונוס}

\noindent{הציגו את מספר התיקון, שם הלקוח ותאריך הקבלה, של התיקונים ש\textbf{דחיפותם אינה "רגיל"} \textbf{ו}שהתקבלו \textbf{לפני \lr{1.7.2026}} \textendash{} ממוינים מ\textbf{התאריך המאוחר לתאריך המוקדם}.}\par
\answerlines{4}

\vspace{8pt}
{\color{inkblue}\hrule height 1pt}
\vspace{4pt}
\begin{center}
{\small\textbf{בהצלחה!} \ \ לפני ההגשה: בדקו שכל שאילתה מסתיימת ב\textbf{נקודה-פסיק}, ושכל ערך טקסט כתוב ב\textbf{גרש בודד}.}
\end{center}

\end{document}
